%=======================================================================
%  2103304 Automatic Control I -- Course Syllabus (Fall 2026)
%  STUDENT VERSION (no CLO/PI mapping)
%  Department of Mechanical Engineering, Chulalongkorn University
%  Compiles with pdfLaTeX.
%=======================================================================
\documentclass[11pt]{article}

\usepackage[T1]{fontenc}
\usepackage[utf8]{inputenc}
\usepackage[margin=1in]{geometry}
\usepackage{array}
\usepackage{booktabs}
\usepackage{tabularx}
\usepackage{enumitem}
\usepackage{xcolor}
\usepackage{colortbl}
\usepackage{titlesec}
\usepackage{mathptmx}      % Times-like serif, clean academic look
\usepackage{fancyhdr}
\usepackage[hidelinks]{hyperref}

% ---- Document version / date (edit here) --------------------------
\newcommand{\docversion}{1.0}
\newcommand{\docdate}{1 August 2026}

% ---- Brand colors (Chula ME) --------------------------------------
\definecolor{brandred}{HTML}{8B2332}
\definecolor{brandred20}{HTML}{ECD7D9}
\definecolor{brandgray}{HTML}{6F6D6E}
\definecolor{boxbg}{HTML}{F7F2F3}

% ---- Section styling ----------------------------------------------
\titleformat{\section}
  {\color{brandred}\normalfont\Large\bfseries}{}{0pt}{}[{\color{brandred!40}\titlerule[1pt]}]
\titlespacing*{\section}{0pt}{16pt}{8pt}

\titleformat{\subsection}
  {\color{brandred}\normalfont\large\bfseries}{}{0pt}{}
\titlespacing*{\subsection}{0pt}{12pt}{4pt}

\setlist[enumerate]{leftmargin=1.5em, itemsep=4pt, topsep=4pt}
\setlist[itemize]{leftmargin=1.4em, itemsep=4pt, topsep=4pt}

\newcolumntype{L}[1]{>{\raggedright\arraybackslash}p{#1}}
\newcolumntype{C}[1]{>{\centering\arraybackslash}p{#1}}
\newcolumntype{R}[1]{>{\raggedleft\arraybackslash}p{#1}}

\renewcommand{\arraystretch}{1.3}
\setlength{\parindent}{0pt}
\setlength{\parskip}{6pt}
\linespread{1.06}\selectfont

% Absorb occasional tight lines by micro-stretching instead of overflowing
\setlength{\emergencystretch}{1em}

% Light horizontal rule between weeks in the schedule
\newcommand{\weekline}{\arrayrulecolor{black!18}\hline\arrayrulecolor{black}}

% Shaded callout box spanning the text width
\newcommand{\callout}[1]{%
  \par\vspace{4pt}%
  \noindent\colorbox{boxbg}{%
    \begin{minipage}{\dimexpr\textwidth-2\fboxsep\relax}
      \vspace{4pt}#1\vspace{4pt}
    \end{minipage}}%
  \par\vspace{6pt}}

% ---- Running header / footer on every page ------------------------
\pagestyle{fancy}
\fancyhf{}
\renewcommand{\headrulewidth}{0.4pt}
\renewcommand{\footrulewidth}{0pt}
\fancyhead[L]{\small\color{brandgray}2103304 Automatic Control I: Course Syllabus (Fall 2026)}
\fancyhead[R]{\small\color{brandgray}Chawin Ophaswongse, Ph.D.}
\fancyfoot[L]{\small\color{brandgray}Version \docversion}
\fancyfoot[C]{\small\color{brandgray}\thepage}
\fancyfoot[R]{\small\color{brandgray}\docdate}
\setlength{\headheight}{14pt}
\setlength{\headsep}{18pt}

\begin{document}

%=======================================================================
%  Title block
%=======================================================================
\begin{center}
    {\LARGE\bfseries Course Syllabus}\\[6pt]
    {\large Department of Mechanical Engineering}\\
    {\large Chulalongkorn University}\\[4pt]
    {\normalsize Academic Year: Fall 2026 (First Semester)}
\end{center}

%=======================================================================
\section*{\color{brandred}Course Information}
\vspace{-2pt}
\begin{tabular}{@{}>{\bfseries}L{5.0cm} L{11.0cm}@{}}
    Course Number   & 2103304 \\
    Course Title    & Automatic Control I \\
    Credits         & 3 (3--0--9) \\
    Semester        & Fall 2026 (First Semester) \\
    Instructor      & Chawin Ophaswongse, Ph.D. \\
    Co-Instructor   & Assoc.\ Prof.\ Thitima Jintanawan, Ph.D. \\
%     Teaching Assistants (TAs) & \dots\\
    Class Time      & Monday, 13:00--16:00 \\
    Location        & Room 409, Building 3, Faculty of Engineering \\
%     Office Hours    & TBD, Hans Bunthli Building, 4th Floor, Room 407 \\
%     Prerequisite    & \dots \\
\end{tabular}

%=======================================================================
\section{Course Description}

\textbf{Official course description.}
\callout{\itshape
Introduction to control systems; mathematical models of systems; feedback control
system characteristics; the performance of feedback control systems; the stability of
linear feedback systems; essential principles of feedback; the root locus method;
frequency response method; stability of the frequency domain, time-domain analysis of
control systems: the design and compensation of feedback control system.}

\textbf{Scope of the course.}
Control engineering is concerned with making a physical system behave in a prescribed
manner, typically faster, more stable, and more accurate than it would be without
control. This course develops the complete workflow, from a physical system to a
validated controller:

\begin{enumerate}
    \item \textbf{Modeling.} Represent a dynamic system of any physical domain, whether
          mechanical, electrical, electro-mechanical, thermal, fluid, or a combination of
          these, by differential equations, transfer functions, and state-space models,
          derived either from first principles or fitted to measured data
          (\emph{system identification}).
    \item \textbf{Analysis.} Determine how the system responds in the time and frequency
          domains, including transient response, steady-state error, and stability.
    \item \textbf{Design.} Select and tune a controller (P, PI, PD, PID, lead/lag, or
          full-state feedback) so that the closed-loop system satisfies specifications
          such as rise time, overshoot, and settling time.
\end{enumerate}

\textbf{Software and hardware.}
All techniques are practiced in \textbf{MATLAB/Simulink} with the Control System
Toolbox. The \emph{Servo Mini-labs} and the \emph{Final Project} extend this work to a
physical servo plant (a brushed DC motor), which students identify, simulate, control,
and tune. All interaction with the hardware takes place through MATLAB/Simulink;
\textbf{no low-level or embedded programming is required}, so that coursework remains
focused on control-engineering decisions rather than firmware implementation.

%=======================================================================
\section{Course Objectives}
Upon successful completion of this course, students will be able to:

\begin{itemize}
    \item \textbf{Derive the governing equations} of a mechanical or electro-mechanical
          system and establish a systematic procedure for analyzing a real engineering
          problem.
    \item \textbf{Construct engineering models} (free-body diagrams, kinematics and
          kinetics diagrams, and control block diagrams) and translate them into
          differential-equation, transfer-function, and state-space representations.
    \item \textbf{Analyze system behavior} in the time and frequency domains, including
          transient response, steady-state error, and stability.
    \item \textbf{Design and tune feedback controllers} (P/PI/PD/PID, lead/lag, and
          full-state feedback) to meet performance specifications, using root-locus and
          frequency-response methods.
    \item \textbf{Apply MATLAB/Simulink} to identify, model, simulate, analyze, and
          design control systems, and validate a controller on a physical servo plant.
\end{itemize}

%=======================================================================
\section{References}
\textbf{Primary}
\begin{itemize}
      \item Nise, N. S. (2011). \emph{Control Systems Engineering} (6th ed.). Wiley.
      % \\\hfill\textit{(Beginner-friendly)}
      \item Lecture and Handout Notes (provided on MyCourseVille)
\end{itemize}
\textbf{Supplementary}
\begin{itemize}
    \item Ogata, K. (2010). \emph{Modern Control Engineering} (5th ed.). Prentice Hall.
          \\\hfill\textit{(State space; MATLAB; more mathematically rigorous)}
    \item Dorf, R. C., \& Bishop, R. H. (2017). \emph{Modern Control Systems} (13th ed.). Pearson.
          \\\hfill\textit{(Design emphasis; rich in illustrative examples)}
\end{itemize}

%=======================================================================
\newpage
\section{Course Schedule}
\begin{center}
\scriptsize
\setlength{\tabcolsep}{4pt}
\begin{tabularx}{\textwidth}{@{}C{0.6cm} C{1.3cm} L{2.9cm} X C{0.7cm} L{3.7cm}@{}}
\toprule
\textbf{Wk} & \textbf{Date} & \textbf{Topic} & \textbf{Key Content} & \textbf{HW} & \textbf{Servo Mini-lab} \\
\midrule
1  & 3 Aug  & Intro \& Math Primer   & Open-/closed-loop $\cdot$ ODEs $\cdot$ LTI $\cdot$ Laplace Transform    & HW1 & \\
\weekline
% 2  & 10 Aug & System Modeling I      & LTI $\cdot$ transfer functions $\cdot$ elec./mech.  & & Lab 1: Intro \& Setup \\
2  & 10 Aug & System Modeling I      & Transfer functions $\cdot$ Electrical/mechanical systems  & & Lab 1: Intro \& Setup \\
\weekline
3  & 17 Aug & System Modeling II     & State space $\cdot$ Nonlinearities $\cdot$ Linearization & HW2 & \\
\weekline
4  & 24 Aug & Time Response          & 1st/2nd-order $\cdot$ Transient-response specifications     & HW3 & Lab 2: Time Response \\
\weekline
5  & 31 Aug & Block Diagrams         & Diagram Reduction $\cdot$ Equivalent forms $\cdot$ Feedback & & \\
\weekline
6  & 7 Sep  & Stability              & Routh--Hurwitz $\cdot$ Gain range for stability   & HW4 & Lab 3: P-Control \& Stability \\
\weekline
7  & 14 Sep & Midterm Review         & Review of Weeks 1--6                               & & \\
\midrule
\multicolumn{6}{@{}l}{\textit{Midterm Examination: Wednesday, 23 September 2026, 13:00--16:00}} \\
\midrule
8  & 28 Sep & Steady-State Errors    & System type $\cdot$ Gain design $\cdot$ Disturbances & HW5 & \\
\weekline
9  & 5 Oct  & Root Locus I           & Sketching rules $\cdot$ Transient response via gain & HW6 & \\
\weekline
10 & 12 Oct & Root Locus II          & PI/PD/PID $\cdot$ Lead/lag compensation            & & Lab 4: Root Locus \\
\weekline
11 & 19 Oct & Frequency Response I    & Asymptotic Bode plots $\cdot$ Bandwidth           & HW7 & \\
\weekline
12 & 26 Oct & Frequency Response II   & Nyquist \& margins $\cdot$ Lead/lag $\cdot$ PID     & & Lab 5: Frequency Response \\
\weekline
13 & 2 Nov  & Design via State Space & Controllability $\cdot$ Full-state feedback         & HW8 & Lab 6: Pole Placement \\
\weekline
14 & 9 Nov  & Project Demo Day       & Group presentations                                & & \\
\weekline
15 & 16 Nov & Final Review           & Review of all topics                              & & \\
\midrule
\multicolumn{6}{@{}l}{\textit{Final Examination: Monday, 30 November 2026, 08:30--11:30}} \\
\bottomrule
\end{tabularx}
\end{center}

%=======================================================================
\section{Course Policies}
\begin{itemize}
    \item \textbf{Communication.} Announcements and course materials are posted on
          MyCourseVille. %; day-to-day discussion takes place on the course Discord.
    \item \textbf{Attendance.} Attendance at mini-lab sessions is \textbf{mandatory and graded}.
    \item \textbf{Mini-labs.} Pre-lab handouts are released one week before each session
          and constitute required reading. Lab reports (groups of up to four) are due \textbf{one
          week} after the session.
    \item \textbf{Assignments.} Homework is due \textbf{ten days} after it is assigned unless
          otherwise specified. Late submissions are penalized according to the stated
          policy.
    \item \textbf{Academic integrity.} Collaboration is encouraged; however, all
          submitted work must be the individual's or the group's own. Plagiarism and
          examination misconduct are handled under Faculty of Engineering regulations.
    \item \textbf{Software.} MATLAB/Simulink and the Control System Toolbox must be
          installed before Week 2.
\end{itemize}

\newpage
%=======================================================================
\section{Assessment}
\begin{center}
\begin{tabular}{@{}L{9cm} C{2cm}@{}}
\toprule
\textbf{Component} & \textbf{Weight} \\
\midrule
Assignments (Homework)      & 10\% \\
Servo Mini-labs             & 10\% \\
Midterm Examination         & 30\% \\
Final Project               & 15\% \\
Final Examination           & 35\% \\
\midrule
\textbf{Total}              & \textbf{100\%} \\
\bottomrule
\end{tabular}
\end{center}


%=======================================================================
% \newpage
\section{Servo Mini-labs}
Six 45-minute hands-on sessions are conducted at the end of the lectures. Students
identify a plant model experimentally (system identification), model its components
(actuator, sensor, and controller), and then design and tune controllers in
\textbf{MATLAB/Simulink} for validation on a real-time servo setup consisting of a
brushed DC motor with an interchangeable inertia-disk or pendulum load. From Lab 2
onward, each session compares the no-load and inertia-disk conditions, so that the
effect of an added load on the same system can be observed directly.

\callout{\textbf{Requirements.} A laptop with MATLAB/Simulink and the Control System
Toolbox is required. Sessions are mandatory and graded. All interaction with the
hardware takes place through the Simulink control interface;
\emph{no low-level or embedded programming is required}.}

% \emph{Tentative session plan (subject to adjustment):}
\begin{itemize}
    \item \textbf{Lab 1: Introduction \& Setup (Wk 2).} Connect the DC-motor and
          inertia-disk plant to Simulink; become familiar with the control interface,
          encoder sensing, and actuation; conduct open-loop step tests and an initial
          system identification of the motor parameters.
    \item \textbf{Lab 2: Time Response (Wk 4).} Measure open- and closed-loop step
          responses; extract first- and second-order parameters together with the
          transient specifications (rise time, overshoot, and settling time) and compare
          them against the identified model, with and without the inertia disk.
    \item \textbf{Lab 3: P-Control \& Stability (Wk 6).} Implement proportional position
          and speed control; observe how the gain shapes the transient response and the
          onset of instability, and how the added inertia changes both.
    \item \textbf{Lab 4: Root Locus (Wk 10).} Design the controller gain (P/PD/PID) by
          root locus to meet transient specifications and validate the design on the
          plant under both load conditions.
    \item \textbf{Lab 5: Frequency Response (Wk 12).} Measure the frequency response of
          the plant (Bode) with and without the inertia disk; evaluate the gain and phase
          margins and tune a lead/lag or PID controller.
    \item \textbf{Lab 6: Pole Placement (Wk 13).} Design full-state feedback by pole
          placement for the DC motor with a load and validate the closed-loop
          response.
\end{itemize}

%=======================================================================
\newpage
\section{Final Project}
Beginning after the midterm, students work in groups of 3--4 to carry out the
complete workflow on an assigned electro-mechanical plant: performing system
identification to obtain the plant model, constructing a \textbf{Simulink simulation
model} and validating it against the measured hardware response, and then designing,
tuning, and demonstrating a controller that meets the target performance specifications.
The comparison between simulation and hardware is a central theme; students are expected
to explain and reconcile the discrepancies they observe.

\emph{Tentative milestones (subject to adjustment):}
\begin{itemize}
    \item \textbf{Milestone 1: System Identification \& Simulation Model (Week 9,
          5 Oct).} Conduct experiments to identify the plant parameters; construct a
          Simulink simulation model and validate it against the measured open-loop
          hardware response.
    \item \textbf{Milestone 2: Controller Design in Simulation (Week 11, 19 Oct).}
          Specify the target performance, design and tune the controller in simulation,
          and verify that the specifications are met on the simulation model.
    \item \textbf{Milestone 3: Hardware Implementation \& Tuning (Week 13, 2 Nov).}
          Deploy the controller on the physical plant, retune as required, and document
          the discrepancies between simulation and hardware.
    \item \textbf{Final Deliverable: Demo Day \& Report (Week 14, 9 Nov).} Present the
          final closed-loop performance at the demo-day and submit the final
          project report.
\end{itemize}

\vspace{6pt}
{\small\color{brandgray}\textit{This syllabus is subject to minor adjustment;
any changes will be announced in class and on MyCourseVille.}}

\end{document}
